\chapter{Structure interne d'un système de fichiers}
	\section{Présentation générale}
	La structure interne dépend du système de fichiers utilisé. Dans le cas d'un système ext2 ou ext3, plusieurs éléments fondamentaux interviennent.

	\section{Le superbloc}
Le superbloc contient les informations essentielles relatives au système de fichiers. Il constitue l'élément central permettant son fonctionnement.

	\section{La table des inodes}
La table des inodes contient les descripteurs des fichiers. Chaque fichier est identifié de manière unique par un numéro d'inode.

	\subsection{Les répertoires}
Les répertoires établissent la correspondance entre le nom d'un fichier et son numéro d'inode.

	\subsection{Les fichiers}
Les fichiers sont identifiés par leur inode et sont considérés comme des suites non ordonnées d'octets.
